\documentclass[12pt]{report}
\usepackage[spanish]{babel}
\usepackage[utf8]{inputenc}
\usepackage[T1]{fontenc}
\usepackage{fancyhdr}
\usepackage{graphicx}
\usepackage{geometry}
\geometry{a4paper, margin=1in}
\usepackage{titlesec}
\usepackage{xcolor}
\usepackage{cite}
\usepackage{parskip}
\usepackage{tocloft}
\usepackage{amsmath}
\usepackage{booktabs}
\usepackage{url}
\renewcommand{\cftsecfont}{\normalfont\MakeUppercase}
\renewcommand{\cftsecpagefont}{\normalfont}
\setlength{\cftsecindent}{20pt}

\definecolor{highlight}{RGB}{255, 255, 0}
\definecolor{azulito}{RGB}{85, 142, 213}

\titleformat{\section}
  {\normalfont\Large\bfseries\MakeUppercase}{}{0pt}{}

\addto\captionsspanish{\renewcommand{\contentsname}{Contenido}}

\renewcommand{\thesection}{}
\renewcommand{\thesubsection}{}


\begin{document}


\thispagestyle{empty}
\vspace*{2cm}
\begin{center}
    \noindent\rule{\textwidth}{0.4pt}\\[0.5cm]
    {\Huge\bfseries MODELO PARA EL ANALISIS DE FACTORES QUE INFLUYEN EN LA POPULARIDAD DE LOS VIDEOJUEGOS}\\[0.5cm]
    \noindent\rule{\textwidth}{0.4pt}\\[0.5cm]
    {\large PROYECTO 01}\\[2cm]
    \textbf{ASIGNATURA:} \hspace{1cm} MÉTODOS NÚMERICOS\\[0.2cm]
    
    \textbf{INTEGRANTES:} \hspace{1cm} RICARDO CISNEROS -- JAIR PAEZ -- HUGO COYAGO\\[2cm]
    
    \textbf{FECHA DE ENTREGA:} \hspace{0.5cm} 21/06/2026
\end{center}
\newpage


\tableofcontents
\thispagestyle{empty}
\newpage

\pagenumbering{arabic}
\setcounter{page}{1}

\section*{Descripción del Proyecto}
\addcontentsline{toc}{section}{DESCRIPCIÓN DEL PROYECTO}

El proyecto tiene la finalidad de permitir el análisis de los distintos factores que influyen en la popularidad de los videojuegos lanzados en la plataforma Steam. Para ello se utilizarán métodos numéricos y análisis de datos sobre una base de datos que contiene la información cruda relacionada con las características de los juegos, métricas de rating y datos del rendimiento comercial.

La problemática de la vida real que se aborda es la competitividad y el desafío de la rentabilización de nuevos lanzamientos en la industria de los videojuegos. Se busca comprender qué variables tienen un mayor impacto en la popularidad de un videojuego, medida por el número de propietarios (ventas estimadas). Con esto se busca indicar los patrones de comportamiento y tendencias de aceptación en el mercado, ayudando así a desarrolladores y publicadores a tomar decisiones informadas para el éxito de sus productos.

Mediante técnicas de aproximación y ajustes de curvas se construyó un modelo numérico que permite analizar el comportamiento de las variables clave y determinar su relevancia en el crecimiento de popularidad.


\section*{Objetivo del Proyecto}
\addcontentsline{toc}{section}{OBJETIVO DEL PROYECTO}

Desarrollar un modelo numérico basado en datos reales de varios videojuegos de la plataforma Steam que permita estimar e identificar los factores que influyen en la popularidad de un videojuego. Empleando métodos numéricos y técnicas de aproximación matemática sobre un conjunto masivo de datos, se busca construir funciones que limpien, normalicen y estructuren la información para resolver un sistema de ecuaciones que indique qué variables (como reseñas, precios, idiomas y características) contribuyen con mayor fuerza al incremento de ventas o propietarios estimados de un videojuego.


\section*{Descripción del Dataset}
\addcontentsline{toc}{section}{DESCRIPCIÓN DEL DATASET}

Para el desarrollo del proyecto se utilizará el dataset “Steam Catalog Insights (October 2024)"\cite{newbie_steam_insights}. Este dataset contiene información masiva relacionada con los videojuegos publicados en la plataforma Steam.

Esta información original se encontraba organizada en múltiples tablas .CSV independientes (Games, Reviews, Steamspy\_insights, Genres, Tags, Categories, Promotional). 

El dataset crudo contaba con más de 140 000 registros en la tabla principal de juegos. Tras el proceso analítico de unificación, integración mediante llaves compartidas (\textit{app\_id}) y la eliminación de datos nulos o inconsistentes, se obtuvo un conjunto de datos consolidado y depurado de 79 215 registros únicos. Estas variables proveen la información técnica, comunitaria y financiera necesaria para la experimentación matemática. La cantidad masiva de datos obtenidos es la adecuada para la alimentación del modelo numérico (70\% de los datos para entrenamiento), la validación de sus ecuaciones (15\% de los datos) y las pruebas predictivas finales (15\% de los datos).


\section*{Variables del Dataset que se usaron}
\addcontentsline{toc}{section}{VARIABLES DEL DATASET QUE USARON}

Para el desarrollo del proyecto se seleccionaron columnas clave estrictamente relacionadas con la aceptación e interacción de los usuarios dentro de la plataforma. La elección de variables como el precio, la cantidad de idiomas disponibles y el año de lanzamiento no es arbitraria. Grm y Šubelj (2021)\cite{grm2021predicting} demostraron que estas características del juego son factores relevantes para estimar tanto el número de jugadores activos como el volumen de ventas durante el ciclo de vida del videojuego.

Las columnas clave explícitamente usadas en el análisis numérico son las siguientes:

\begin{table}[h]
\centering
\begin{tabular}{|p{3cm}|p{8cm}|p{3cm}|}
\hline
\textbf{Variable} & \textbf{Descripción} & \textbf{Presentación} \\ \hline
estimated\_owners & Número estimado de propietarios del videojuego & float \\ \hline
approval\_ratio & Proporción de aprobación calculada a partir de reseñas positivas y negativas & 0.0 a 1.0 \\ \hline
is\_free & Indicador binario de si el videojuego es gratuito (1) o de pago (0) & Binario (1 gratis y 0 pago) \\ \hline
reviews\_total & Volumen acumulado de críticas de la comunidad. & float \\ \hline
price\_usd & Precio del videojuego en dólares normalizado en dólares continuos. & Centavos/100 \\ \hline
category\_count & Muestra la Cantidad total de modos de juego y características técnicas que ofrece el videojuego. & float \\ \hline
language\_count & Cantidad total de idiomas de locación y traducción disponibles en el videojuego. & float \\ \hline
release\_year & Año cronológico del lanzamiento extraído de la fecha de publicación oficial. & 1997 a 2024 \\ \hline
\end{tabular}
\end{table}

Justificación: Estas variables fueron seleccionadas debido a que poseen una relación directa comprobable con el comportamiento de los usuarios y el desempeño comercial. Al usar el total de reseñas y la aprobación como variables principales, en lugar de puntuaciones de críticos especializados, nos respaldamos en la literatura contemporánea. Robinson (2020)\cite{robinson2020predicting} demostró que las opiniones directas de los jugadores y la visibilidad comunitaria son métricas predictivas de ventas acumuladas mucho más estables que las evaluaciones tradicionales de prensa. Finalmente, evitan la colinealidad fuerte y los sesgos que presentan variables más subjetivas.


\section*{Resolución}
\addcontentsline{toc}{section}{RESOLUCIÓN}

\subsection*{Fases que propusieron en el primer bimestre}
Para llevar a cabo el proyecto estructurado, se propusieron inicialmente las siguientes fases operativas:
\begin{enumerate}
    \item \textbf{Recolección e integración de datos:} Consolidación de múltiples tablas CSV separadas en un único repositorio integrado utilizando un identificador maestro.
    \item \textbf{Limpieza y preparación estadística:} Depuración de nulos, eliminación de duplicados, corrección de tipos de dato y normalización estadística de variables.
    \item \textbf{Planteamiento del Modelo Matemático:} Construcción algorítmica de Ecuaciones Normales por el principio de Mínimos Cuadrados para crear una matriz de diseño.
    \item \textbf{Implementación Computacional:} Desarrollo lógico y desde cero de métodos numéricos exactos (Solucionadores Directos) en el lenguaje Python.
    \item \textbf{Evaluación y Validación:} Pruebas de rendimiento del modelo sobre datos desconocidos con la aplicación de métricas de precisión (MAE, RMSE, R²).
\end{enumerate}

\subsection*{Etapa del análisis (incluye fase de limpieza de datos)}
Durante esta etapa se realizó la limpieza y adecuación del dataset masivo. Se eliminaron registros duplicados y se descartaron valores nulos o vacíos en las variables clave. Se aplicaron validaciones que aseguraran la consistencia física de los datos (por ejemplo: que los precios no sean negativos, el año de lanzamiento estuviera entre 1997 y 2024, la aprobación estuviese contenida en [0.0, 1.0] y los propietarios fuesen mayores a cero).

Debido a los fuertes cambios de escala (precios en decenas contra reseñas en millones), se ejecutó un proceso de escalamiento estandarizado. Tanto las siete variables independientes como la variable dependiente (estimación de propietarios) fueron normalizadas con el método de Z-score, transformando todos los datos a un dominio centrado en 0 con desviación estándar de 1, evitando así la inestabilidad numérica.

\subsection*{Métodos numéricos utilizados dentro de la resolución}
Se utilizó el método fundamental de Ecuaciones Normales de Mínimos Cuadrados multivariable para construir una función de aproximación que represente las tendencias de los 55,450 videojuegos utilizados para entrenamiento. El modelo se plantea como:
$$ y = a_0 + a_1x + a_2x^2 + \dots + a_nx^n $$
Donde $y$ es el valor normalizado de los propietarios y $x$ es el vector de predictores.

A nivel matricial, esto dio como resultado la reducción del problema a un sistema de ecuaciones lineales compacto de dimensiones $8 \times 8$:
$$ A c = b $$

Para resolver de manera exacta y con alta estabilidad dicho sistema matricial, se implementaron dos métodos numéricos directos:
\begin{itemize}
    \item \textbf{Eliminación Gaussiana con Pivoteo Parcial:} Para transformar la matriz a una forma triangular superior sin divisiones por ceros y sustituir hacia atrás.
    \item \textbf{Descomposición LU (Algoritmo de Doolittle):} Factorizando la matriz original en dos matrices triangulares complementarias ($L$ y $U$) posibilitando su resolución ágil mediante sustitución progresiva y regresiva.
\end{itemize}

\subsection*{Explicación breve de las funciones implementadas}
Para la estructura de datos, se usó \textbf{DataFrames de Pandas} para el manejo y filtrado veloz de tablas y CSV, y \textbf{Arrays de NumPy} para crear matrices puras donde se aplicarían los cálculos algebraicos y multiplicaciones matriciales.

\begin{itemize}
    \item \textbf{Carga y unificación:} Funciones que leen los datos, filtran columnas útiles y las combinan haciendo coincidir los `app_id`.
    \item \textbf{Limpieza:} Funciones que utilizan mascaras y condicionales lógicos para filtrar registros que exceden límites naturales, aplicando `dropna()` para registros faltantes.
    \item \textbf{Normalización:} Uso de operaciones vectorizadas matemáticamente puras (restar el promedio y dividir por la desviación estándar) creando nuevas columnas normalizadas.
    \item \textbf{División de datos:} Construcción de matrices de diseño usando `np.hstack` para unir un vector de unos (para el intercepto) con los valores normalizados, y uso de `iloc` para repartir filas al azar entre Train/Val/Test.
    \item \textbf{Solucionadores Gauss y LU:} Estructuras puramente algorítmicas (iteraciones `for` anidadas) que manipulan la posición de arreglos de índices `A[i][j]`.
\end{itemize}

\subsection*{Flujograma de interacción entre funciones implementadas}

\begin{enumerate}
    \item \textbf{Ingesta e Integración (Funciones auxiliares):} El programa inicia cargando librerías clave (\texttt{pandas}, \texttt{numpy}) y declarando funciones de pre-procesamiento como \texttt{analizar\_rango\_dueños()} y \texttt{extraer\_año\_lanzamiento()}. Posteriormente, utiliza la función \texttt{pd.read\_csv()} iterativamente para cargar las bases de datos originales, fusionándolas a través de \texttt{pd.merge()} basándose en la llave \textit{app\_id} y calculando matemáticamente las nuevas variables continuas (precio en dólares y ratio de aprobación).
    \item \textbf{Ciclos de Limpieza y Condicionales (Rombos lógicos):} La base de datos unificada se pasa por sentencias de validación apoyadas en la función \texttt{dropna()} para purgar filas incompletas. Seguidamente, se implementan estructuras condicionales compuestas (operadores lógicos sobre el DataFrame) para evaluar si el precio es $\geq 0$, si el año está entre 1997-2024 y si el ratio de aprobación $\in [0, 1]$. Los registros que no cumplen estas sentencias lógicas son excluidos del flujo de ejecución, depurando la muestra.
    \item \textbf{Transformación Z-score y \textit{Slicing}:} A los datos validados se les ejecuta el escalamiento normalizado Z-score mediante aritmética de arreglos. Usando la función de Pandas \texttt{sample(frac=1, random\_state=87)}, el conjunto se baraja aleatoriamente. Posteriormente, se realiza un particionamiento de las filas usando índices espaciales (\texttt{.iloc}) para dividir los subconjuntos en variables independientes para entrenamiento (70\%), validación (15\%) y pruebas (15\%).
    \item \textbf{Construcción Matricial y Solucionadores (Bucles Anidados):} Utilizando \texttt{np.dot()} de NumPy, el programa ejecuta el producto cruzado transpuesto para conformar algebraicamente la Matriz de Ecuaciones Normales ($A = X^T X$) y su respectivo vector ($b = X^T y$). A continuación, se invocan a las funciones personalizadas \texttt{solucion\_directa\_gauss()} y \texttt{solucion\_directa\_LU()}. Estas funciones dependen de \textbf{múltiples bucles \texttt{for} y \texttt{while} anidados} que iteran minuciosamente sobre los índices matriciales bidimensionales \texttt{[i][j]} para ejecutar operaciones elementales de fila (pivoteo, eliminación progresiva y sustitución regresiva).
    \item \textbf{Métricas de Evaluación y Diagnóstico de Anomalías:} Con el vector solución devuelto por los solucionadores, se evalúa el rendimiento general mediante operaciones vectorizadas que calculan \texttt{R²}, \texttt{RMSE} y \texttt{MAE}. Finalmente, el programa genera un arreglo de residuos estandarizados y aplica una compuerta iterativa sobre todo el conjunto de prueba: mediante la sentencia lógica \texttt{if |z\_score| > 3}, el código clasifica los registros como valores atípicos (anomalías), aislándolos en memoria para imprimir un "Top 5" de los peores errores de predicción, y concluye midiendo tiempos de procesamiento con la librería \texttt{time} y generando las gráficas comparativas.
\end{enumerate}


\subsection*{Link de repositorio}
Todo el contenido del proyecto se encuentra publicado en el siguiente repositorio de GitHub:

\url{https://github.com/Jair-paez/Grupo5-MetodosNumericos}

\subsection*{Estadísticas de contribuciones del repositorio}
El trabajo se estructuró a través del sistema de control de versiones Git, garantizando que el equipo pudiera desarrollar la matemática, el reporte y la limpieza de datos paralelamente. El registro oficial de *commits* arroja las siguientes estadísticas de contribución:



\subsection*{Resultado final, tiempos de ejecución y testeo de simulación}

Al resolver el sistema $8 \times 8$ resultante, ambos solucionadores (Gauss y LU) arrojaron resultados matemáticamente idénticos, validando la robustez de la implementación. 

Dado que el problema gigantesco se comprimió a través de Ecuaciones Normales a una simple matriz de 8 filas y columnas, los tiempos de ejecución de la simulación o resolución del modelo lineal fueron cpracticamente instantáneos (fracciones invisibles de segundo). Esto corrobora la enorme eficiencia matemática de aplicar métodos directos tras una buena reducción matricial.

El análisis de los pesos de la ecuación reveló que el coeficiente de mayor influencia es el total de reseñas normalizado ($a_4 = 0.7988$), seguido del conteo de idiomas.

\textbf{Validación contra Datos Desconocidos:}
\begin{itemize}
    \item Se probaron los coeficientes contra el dataset de \textbf{validación} (15\% de datos), obteniéndose un $R^2$ de \textbf{78.43\%}, con errores ínfimos (MAE de 0.0673).
    \item En la simulación final, se probó el modelo frente al conjunto estricto de \textbf{testeo} (15\% restante), logrando un $R^2$ de \textbf{76.95\%}.
\end{itemize}

A partir del residuo estandarizado se detectaron 304 valores atípicos (0.38\% de los datos). Estos representan una proporción pequeña del conjunto de datos, ya que la mayor parte de los registros mantiene errores controlados. Los valores atípicos identificados corresponden a los títulos que se convierten en mega-éxitos o fenómenos de la cultura popular masiva; evidenciando que el comportamiento matemático generalizado del mercado es estable, pero requiere consideración para anomalías o tendencias de viralidad (factores externos).


\section*{Conclusiones}
\addcontentsline{toc}{section}{CONCLUSIONES}
\begin{itemize}
    \item El análisis de una cantidad grande datos permitió identificar matemáticamente que la popularidad de un videojuego no está ligada a solo un factor, sino de la combinación ponderada de múltiples variables relacionadas con la interacción de los usuarios.
    \item El uso de técnicas metódicas de limpieza, integración y normalización Z-score garantizó la convergencia sin sesgos durante la aplicación del método de mínimos cuadrados, siendo vital el tratamiento adecuado de matrices para prevenir la pérdida de estabilidad numérica.
    \item La ejecución paralela de la Eliminación Gaussiana y la Descomposición LU sirvió para demostrar la robostuza del modelo.
\end{itemize}

\section*{Recomendaciones}
\addcontentsline{toc}{section}{RECOMENDACIONES}
\begin{itemize}
    \item Usar datasets provenientes de fuentes confiables, de preferencia proporcionados por empresas o repositorios de prestigio (como las bases estandarizadas de GitHub utilizadas en esta entrega) para asegurar la validez real de las conclusiones algorítmicas extraídas.
    \item Realizar una interpretación minuciosa de correlaciones o colinealidad antes de seleccionar las variables predictoras, puesto que el uso innecesario de variables correlacionadas engrosa artificialmente las matrices incrementando el tiempo de procesamiento y el margen de error.
    \item Emplear un único modelo de normalizacion, para evitar comparaciones que puedan generar sesgos o predicciones erroneas.
\end{itemize}

\vspace{0.5cm}

%\newpage
\renewcommand{\bibname}{\MakeUppercase{REFERENCIAS}}
\bibliographystyle{IEEEtran}
\addcontentsline{toc}{section}{REFERENCIAS}
\bibliography{inFiles/References/references}

\end{document}
